\chapter{Simulation and Robot Control}
\label{ch:robots}
\chapteroverview{Explain how Gazebo assets are resolved, how robot instances are
selected and isolated, and how validated commands cross the ROS--Gazebo
boundary.}{Adding or selecting a robot, diagnosing missing assets or topics,
changing arm/gripper behavior, or evaluating simulation-specific limitations.}

\section{Chapter Source Map}

\small
\begin{longtable}{L{0.23\textwidth} L{0.69\textwidth}}
\caption{Authoritative sources for simulation and robot control}
\label{tab:robot-source-map}\\
\toprule
\tableheader Topic & Authoritative repository path(s) \\
\midrule
\endfirsthead
\toprule
\tableheader Topic & Authoritative repository path(s) \\
\midrule
\endhead
World and model resolution & \file{mfja_3rd_floor_description/CMakeLists.txt}; \file{mfja_3rd_floor_description/worlds/}; \file{mfja_3rd_floor_description/models/}; \file{mfja_robot_control_config/launch/multi_robot_sim.launch.py} \\ \rowseparator
Robot registry and selection & \file{mfja_robot_control_config/config/robots.yaml}; \file{mfja_robot_control_config/config/robots_room_315_only.yaml}; \file{mfja_robot_control_config/launch/multi_robot_sim.launch.py} \\ \rowseparator
Joint command and trajectory control & \file{mfja_robot_control_config/scripts/robot_joint_command.py}; \file{mfja_3rd_floor_description/src/SmoothJointTrajectoryController.cc}; industrial and TIAGo \file{model.sdf} plugin declarations \\ \rowseparator
Gripper configuration and control & \file{mfja_robot_control_config/config/gripper_command_defaults.yaml}; \file{mfja_robot_control_config/launch/gripper_range_config.py}; \file{mfja_robot_control_config/scripts/robot_gripper_command.py}; \file{mfja_3rd_floor_description/src/SymmetricGripperController.cc} \\ \rowseparator
TIAGo frame isolation & \file{mfja_robot_control_config/launch/multi_robot_sim.launch.py}; \file{mfja_3rd_floor_description/models/tiago/model.sdf}; \file{mfja_3rd_floor_description/models/tiago_base/model.sdf}; \file{mfja_3rd_floor_description/urdf/tiago.urdf}; \file{mfja_3rd_floor_description/urdf/tiago_base.urdf} \\ \rowseparator
Industrial collision boundary & \file{mfja_3rd_floor_description/models/kuka_kr6r900sixx/model.sdf}; \file{mfja_3rd_floor_description/models/staubli_tx2_60l/model.sdf}; \file{mfja_3rd_floor_description/models/yaskawa_hc10/model.sdf}; \file{mfja_3rd_floor_description/models/yaskawa_hc10dt/model.sdf} \\
\bottomrule
\end{longtable}
\normalsize

\section{Asset and Model Resolution}

Every maintained reusable Gazebo model normally has this form:

\begin{lstlisting}
mfja_3rd_floor_description/models/<model_name>/
  model.config
  model.sdf
  meshes/...        # where required
\end{lstlisting}

World SDF includes models with \code{model://...} URIs. The launch prepends the
description model directory to Gazebo's model and resource paths while
preserving any pre-existing entries. A missing model/mesh error usually means
the description package was not built/sourced, a URI does not match its model
directory, or a local environment overwrote the resource path.

The main worlds include static room/furniture/cell assets. Robots are not fixed
world composition: they are selected from YAML and spawned after Gazebo starts.
Eight identity-specific Room 315 shuttle models are preloaded below the scene
and are revealed or managed by the rail nodes.

\section{Robot Instances and Runtime Boundary}

\subsection{Robot registry}

The authoritative instance, model, spawn-pose, and profile-default mapping is
kept once in Table~\ref{tab:robot-registry-contract}. Its owning configuration
files are \file{mfja_robot_control_config/config/robots.yaml} and
\file{mfja_robot_control_config/config/robots_room_315_only.yaml}; this chapter
focuses on how those selected instances cross the runtime boundary.

The full-floor registry enables all six listed instances. TIAGo aliases include
\code{tiago}, \code{tiago_with_arm}, and \code{tiago_arm}; base aliases include
\code{tiago_base}, \code{tiago_no_arm}, and \code{tiago_mobile_base}. A mobile
entry may also use variant/type fields or an explicit arm boolean. Names must be
unique because each name becomes a Gazebo entity, ROS namespace, bridge-file
name, and TF prefix.

\subsection{Per-robot data flow}

\begin{longtable}{L{0.17\textwidth} L{0.17\textwidth} L{0.34\textwidth} L{0.24\textwidth}}
\caption{Per-instance robot command and feedback boundaries}
\label{tab:robot-data-flow}\\
\toprule
\tableheader Stage & Direction & Responsibility & Source file(s) \\
\midrule
\endfirsthead
\toprule
\tableheader Stage & Direction & Responsibility & Source file(s) \\
\midrule
\endhead
Validated helper & Operator $\rightarrow$ ROS & Resolve one registry instance, validate the request, and publish only to its namespace & \file{mfja_robot_control_config/scripts/robot_joint_command.py}; \file{mfja_robot_control_config/scripts/robot_gripper_command.py} \\ \rowseparator
Generated bridge & ROS $\leftrightarrow$ Gazebo & Map instance-specific command, state, odometry, and TF topics & \file{mfja_robot_control_config/launch/multi_robot_sim.launch.py} \\ \rowseparator
Model controller & Gazebo transport $\rightarrow$ model & Apply bounded joint, gripper, or mobile-base motion inside Gazebo & \file{mfja_3rd_floor_description/src/SmoothJointTrajectoryController.cc}; \file{mfja_3rd_floor_description/src/SymmetricGripperController.cc}; selected model \file{model.sdf} \\ \rowseparator
Feedback lane & Gazebo $\rightarrow$ ROS & Return joint state, progress, odometry, and TF through the same instance boundary & \file{mfja_robot_control_config/launch/multi_robot_sim.launch.py}; selected URDF/model SDF \\
\bottomrule
\end{longtable}

There is no global broadcast joint-command topic.

\subsection{Canonical robot ROS topics}

Replace \code{<name>} with the registry instance name.

\begin{longtable}{L{0.36\textwidth} L{0.27\textwidth} L{0.27\textwidth}}
\caption{Canonical namespaced ROS topics exposed for robot instances}
\label{tab:robot-topics}\\
\toprule
\tableheader ROS name/type & Direction & Availability \\
\midrule
\endfirsthead
\toprule
\tableheader ROS name/type & Direction & Availability \\
\midrule
\endhead
\rosname{/<name>/joint_trajectory} --- \code{trajectory_msgs/JointTrajectory} & ROS $\rightarrow$ Gazebo & Every robot \\ \rowseparator
\rosname{/<name>/joint_trajectory_progress} --- \code{std_msgs/Float32} & Gazebo $\rightarrow$ ROS & Every robot using the smooth plugin \\ \rowseparator
\rosname{/<name>/joint_states} --- \code{sensor_msgs/JointState} & Gazebo $\rightarrow$ ROS & Every robot \\ \rowseparator
\rosname{/<name>/tf}, \rosname{/<name>/tf_static} & State publisher/Gazebo $\rightarrow$ ROS & Every robot; mobile TF also comes from DiffDrive \\ \rowseparator
\rosname{/<name>/gripper/position_command} --- \code{std_msgs/Float64} & ROS $\rightarrow$ Gazebo & Industrial robots \\ \rowseparator
\rosname{/<name>/cmd_vel} --- \code{geometry_msgs/Twist} & ROS $\rightarrow$ Gazebo & TIAGo variants \\ \rowseparator
\rosname{/<name>/odom} --- \code{nav_msgs/Odometry} & Gazebo $\rightarrow$ ROS & TIAGo variants \\
\bottomrule
\end{longtable}

\section{Industrial Arm Control}

\subsection{Joint command helper}

\file{robot_joint_command.py} resolves the robot profile, validates joint
names/counts/limits, waits for a complete live joint state, and publishes once a
dense quintic trajectory sampled at 100--200 Hz to the selected namespace.
Inspect its live help because supported options are executable authority:

\begin{lstlisting}[language=bash]
ros2 run mfja_robot_control_config robot_joint_command.py --help
ros2 run mfja_robot_control_config robot_joint_command.py --list
\end{lstlisting}

Use its dry-run/validation path before a new pose. The helper's dense message is
interpolated as a multi-point trajectory by the smooth controller, which retimes
motion when necessary to remain within configured joint velocity limits. The
controller rejects unknown/duplicate joints, non-finite values, mismatched point
sizes, non-monotonic times, and positions outside SDF limits.

\subsection{Trajectory controller semantics}

\file{SmoothJointTrajectoryController.cc} is deterministic and simulation-time
driven. It uses joint position/velocity reset components on every step, then
rewrites public joint state after physics so gravity/solver residue does not
leak into feedback.

\begin{itemize}
  \item One-point trajectories are goals: current state to target through a
        quintic minimum-jerk blend.
  \item Multi-point trajectories are timestamped samples and interpolate
        linearly between points after velocity-limit retiming.
  \item A newer valid message replaces the active trajectory.
  \item Commands received while paused are retained; progress starts when
        simulation time advances.
  \item Progress is published on the Gazebo trajectory topic plus
        \code{_progress}, at no more than one update per 0.05 simulated seconds.
\end{itemize}

\section{Gripper Configuration and Control}

\file{gripper_command_defaults.yaml} declares, per industrial instance, the
per-jaw position at 0\% and 100\%, plus default open/close percentages. Current
100\% per-jaw travels are 0.030 m (KUKA), 0.010 m (St\"aubli), 0.040 m (HC10),
and 0.010 m (HC10DT).

At launch, \file{gripper_range_config.py} validates the selected robots and
config, then materializes temporary SDF/URDF content with consistent joint and
plugin limits. \file{robot_gripper_command.py} maps an operator percentage or
bounded position into \rosname{/<name>/gripper/position_command}. Inspect before
actuation:

\begin{lstlisting}[language=bash]
ros2 run mfja_robot_control_config robot_gripper_command.py --help
\end{lstlisting}

\begin{safetybox}[Gripper realism]
The symmetric plugin animates two opposite-axis jaws at bounded velocity. It
does not create a grasp constraint, attach a payload, or validate contact force.
A visually enclosed object is not a physically held object.
\end{safetybox}

\section{TIAGo Frame Isolation}

For each mobile instance, the launch creates an instance-specific SDF copy that
rewrites \code{cmd_vel}, odometry, TF, frame, and child-frame names. Gazebo
DiffDrive publishes \code{<name>/odom -> <name>/base_link}; the namespaced state
publisher adds the remaining tree with \code{frame_prefix=<name>/}. Multiple
TIAGo variants can therefore share a graph without colliding on
\code{odom}, \code{base_link}, or downstream frames.

\section{Robot Model Limitations}

Industrial assets are optimized for visual integration and intentionally use
collision simplification; some are collision-free. Visual overlap is not proof
of safe physical clearance. Commands are not ROS 2 control actions and there is
no hardware controller manager in this stack. Before adapting a robot to
hardware, replace the simulation bridge/controller boundary and establish
hardware limits, collision monitoring, e-stop behavior, and vendor-specific
safety integration independently.
